\section{线性椭圆方程}\label{chap:linear-elliptic-equations}

\subsection{定义}\label{sec:elliptic-equation-definitions}

从考察 \eqref{eq:general-linear-differential-expression} 中定义的形如 $L(\bm{x},D)u=0$ 的微分算子开始研究线性椭圆 PDE. 注意, 根据 Definition~\ref{def:higher-order-pde-type}, $m$ 阶微分算子在 $\bm{x}_0$ 处椭圆, 当且仅当
\begin{equation}\label{eq:elliptic-opening-principal-symbol-nonzero}
L^p(\bm{x}_0,\bm{\xi})=\sum_{|\alpha|=m}a_\alpha(\bm{x})\bm{\xi}^\alpha\ne0\quad\text{for every }\bm{\xi}\in\mathbb{R}^n\backslash\{0\}.
\end{equation}
事实上, 可以证明以下结论.

\begin{Lemma}
\label{lem:elliptic-operator-even-order-fixed-sign}
若 $m$ 阶线性偏微分算子 $L$ 在 $\bm{x}_0\in\mathbb{R}^n$, $n>1$, 处椭圆, 则 $m$ 为偶整数 ($m=2k$), 且 $\bm{\xi}\mapsto L^p(\bm{x}_0,\bm{\xi})$ 在 $\bm{\xi}\ne0$ 上只取一种符号.
\end{Lemma}

\begin{Proof}
由定义, $\bm{\xi}\mapsto L^p(\bm{x}_0,\bm{\xi})$ 连续, 且只在 $\bm{\xi}=0$ 处取值 $0$. 假设 $L^p(\bm{x}_0,\bm{\xi}_1)<0$ 且 $L^p(\bm{x}_0,\bm{\xi}_2)>0$, 然后用一条不经过 $0$ 的路径连接 $\bm{\xi}_1$ 和 $\bm{\xi}_2$. 如前所述, $L^p(\bm{x}_0,\bm{\xi})$ 必须沿路径连续变化, 从而取到值 $0$. 这产生矛盾.

现在得到, 对任意 $\bm{\xi}\in\mathbb{R}^n$,\[L^p(\bm{x}_0,\bm{\xi})\quad\text{and}\quad L^p(\bm{x}_0,-\bm{\xi})=(-1)^mL^p(\bm{x}_0,\bm{\xi})\]必须同号. 这说明 $m$ 为偶数.
\end{Proof}

鉴于这一结果, 在本章余下部分采用以下稍加限制的椭圆算子定义.

\begin{Definition}
\label{def:positive-sign-elliptic-and-uniform-elliptic}
设 $\Omega\subseteq\mathbb{R}^n$ 为区域. 若线性偏微分算子
\begin{equation}\label{eq:even-order-elliptic-differential-operator}
L(\bm{x},D)=\sum_{|\alpha|\leq2k}a_\alpha(\bm{x})D^\alpha
\end{equation}
满足
\begin{equation}\label{eq:positive-sign-ellipticity-condition}
(-1)^k\sum_{|\alpha|=2k}a_\alpha(\bm{x})\bm{\xi}^\alpha>0\quad\text{for every }\bm{x}\in\Omega,\ \bm{\xi}\in\mathbb{R}^n\backslash\{0\}.
\end{equation}
则称其在 $\Omega$ 中\textbf{椭圆}. 若存在常数 $\theta>0$, 使得
\begin{equation}\label{eq:uniform-ellipticity-principal-symbol-bound}
(-1)^k\sum_{|\alpha|=2k}a_\alpha(\bm{x})\bm{\xi}^\alpha\geq\theta|\bm{\xi}|^{2k}\quad\text{for every }\bm{x}\in\Omega,\ \bm{\xi}\in\mathbb{R}^n\backslash\{0\}.
\end{equation}
则称 $L$ 在 $\Omega$ 中\textbf{一致椭圆}.
\end{Definition}

\begin{Example}
\label{ex:laplacian-biharmonic-uniform-ellipticity}
读者应当回想 Chapter~\ref{chap:characteristics} 中的计算, 其中表明 Laplacian 的负值 $-\Delta$ (阶数为 $2$) 和双调和算子 $\Delta^2$ (阶数为 $4$) 均以 $\theta=1$ 一致椭圆.
\end{Example}

\begin{Example}
\label{ex:second-order-negative-matrix-uniform-ellipticity}
在 $n$ 维空间中, 形如
\begin{equation}\label{eq:second-order-elliptic-nondivergence-operator}
Lu=a_{ij}(\bm{x})\frac{\partial^2u}{\partial x_i\partial x_j}+b_i(\bm{x})\frac{\partial u}{\partial x_i}+c(\bm{x})u
\end{equation}
的二阶算子在区域 $\Omega$ 上一致椭圆, 只要存在常数 $\theta$, 使得对每个 $\bm{x}\in\Omega$ 均有
\begin{equation}\label{eq:second-order-ellipticity-negative-coefficient-matrix}
\bm{\xi}^T\bm{A}(\bm{x})\bm{\xi}>\theta|\bm{\xi}|^2
\end{equation}
这里 $\bm{A}(\bm{x})$ 为分量为 $-a_{ij}(\bm{x})$ 的 $n\times n$ 矩阵.
\end{Example}

在下面关于存在性与正则性理论的讨论中, 将微分算子写成便于分部积分的形式会很方便.

\begin{Definition}
\label{def:general-divergence-form-operator}
若存在函数 $a_{\sigma\gamma}:\Omega\to\mathbb{R}$, 使得
\begin{equation}\label{eq:general-even-order-divergence-form-definition}
L(\bm{x},D)u=\sum_{0\leq|\sigma|,|\gamma|\leq k}(-1)^{|\sigma|}D^\sigma(a_{\sigma\gamma}(\bm{x})D^\gamma u).
\end{equation}
则称该算子具有\textbf{散度形式}.
\end{Definition}

\begin{Remark}
\label{rem:divergence-form-ellipticity-symbol}
注意, 散度形式的算子椭圆, 当且仅当
\begin{equation}\label{eq:divergence-form-ellipticity-symbol-condition}
\sum_{|\sigma|,|\gamma|=k}\bm{\xi}^\sigma a_{\sigma\gamma}(\bm{x})\bm{\xi}^\gamma>0\quad\text{for every }\bm{x}\in\Omega,\ \bm{\xi}\in\mathbb{R}^n\backslash\{0\},
\end{equation}
它一致椭圆, 当且仅当存在 $\theta>0$, 使得
\begin{equation}\label{eq:divergence-form-uniform-ellipticity-condition}
\sum_{|\sigma|,|\gamma|=k}\bm{\xi}^\sigma a_{\sigma\gamma}(\bm{x})\bm{\xi}^\gamma>\theta|\bm{\xi}|^{2k}\quad\text{for every }\bm{x}\in\Omega,\ \bm{\xi}\in\mathbb{R}^n\backslash\{0\}.
\end{equation}

\end{Remark}
